\documentclass[11pt]{article}
\usepackage[T1]{fontenc}
\usepackage[utf8]{inputenc}
\usepackage{lmodern,microtype,geometry,amsmath,amssymb,amsthm,mathtools,enumitem}
\usepackage[colorlinks=true,linkcolor=blue,urlcolor=blue,citecolor=blue]{hyperref}
\geometry{margin=1in}
\setlength{\emergencystretch}{3em}
\hypersetup{pdftitle={Quadratic-width and linear-depth learning bounds for Bayesian networks with stable populations},pdfauthor={ChatGPT-6 Astra; Claude Fable 5 and Opus 5; Kaarel Hanni; Dmitry Vaintrob}}
\newtheorem{theorem}{Theorem}
\newtheorem{lemma}[theorem]{Lemma}
\newtheorem{corollary}[theorem]{Corollary}
\newtheorem{remark}[theorem]{Remark}
\newcommand{\R}{\mathbb R}
\newcommand{\E}{\mathbb E}
\newcommand{\Prb}{\mathbb P}
\newcommand{\norm}[1]{\left\|#1\right\|}
\newcommand{\clip}{\mathrm{clip}}
\title{Quadratic-width and linear-depth learning bounds\protect\\for Bayesian networks with stable populations}
\author{ChatGPT-6 Astra \and Claude Fable 5 and Opus 5 \and Kaarel H\"anni \and Dmitry Vaintrob\thanks{Authors are listed in alphabetical order.}}
\date{Research draft, 29 October 2026}
\begin{document}
\maketitle

\begin{abstract}
We prove sample bounds comparable to the parameter count of a small
teacher, independently of arbitrarily larger learner width, for exact
Bayesian learning with Gaussian tanh priors restricted by stability and
redundancy of neuron populations.
For a Boolean teacher of width $C$ and $L-1$ gate stages, a digital-population
condition gives sufficient samples $O(dC+LC^2)$ at fixed calibration,
accuracy, and confidence, with no hidden logarithm of $C$ or $L$.
The condition requires that averaging neurons preserve predictions, that
the averaged hidden states be nearly Boolean with a fixed relative margin
at every gate, and that the number of populations be calibrated by readout
energy. It is specified without knowing the target or its complexity.
Every eligible Boolean circuit has a positive-Gaussian-mass implementation
satisfying it. A condition allowing analog hidden states instead requires
internal-noise stability to persist after averaging; it gives
$\widetilde O(dC+LC^2)$, with a stronger per-layer version giving
$O(C^2L\log^5(2+C))$ at fixed other resources.
All guarantees hold simultaneously at every learner width $N\ge C$ and
every positive mean-field-or-stronger Gaussian scale, including finite
mean field. These are different structural prior restrictions, not an
improvement for unrestricted networks or for internal-noise stability alone.
The proofs do not establish that ordinary trained networks satisfy the
conditions, and are not Lean formalized.
\end{abstract}

\section{Introduction}

A dense network with width $C$ and depth $L$ has roughly $dC+LC^2$
parameters. A Bayesian learner may have arbitrarily larger width, so counting
its own parameters does not explain learning a target with a small
implementation. A useful statistical goal is a sample bound comparable,
up to logarithms, to the parameter count of the smaller implementation.

We obtain such a bound for Boolean-circuit targets by restricting the prior
in three ways. The network must admit averaging into a limited number of
neuron populations; this averaging must preserve its clipped predictions;
and the averaged network itself must tolerate finite simultaneous internal
disturbances. The population budget is calibrated by the network's own
readout energy and the architecture, not by the unknown target complexity.
The restriction has positive Gaussian probability and admits every Boolean
teacher considered below.

The new stability requirement matters for the depth dependence. A generic
parameter cover propagates rounding errors through all layers and therefore
needs an exponentially fine grid. Here the rounded network is viewed instead
as the unrounded network with an explicitly constructed disturbance after
each activation. Tanh bounds the rounded hidden states, so the total
required disturbance is polynomially small in the grid mesh and depth.
Finite-radius robustness then controls the final clipped output without
multiplying layer norms.

There are three precise versions of the result:
\begin{enumerate}[leftmargin=*]
\item Robustness of the averaged representation on the usual aggregate
noise ball gives the nearly parameter-count bound of
Theorem~\ref{thm:main}.
\item Allowing each layer its full calibrated noise budget gives a bound
literally linear in depth, with only width logarithms
(Theorem~\ref{thm:perlayer}).
\item Saturation and relative gate margins permit a constant-precision
digital code. With a stronger population-energy calibration this gives
literal quadratic width and linear depth
(Theorem~\ref{thm:digital}). The digital and analog conditions are different
support restrictions; neither is asserted to follow from the other.
\end{enumerate}
The first construction counts dense compressed-network parameters. The
last counts constant-radius nets for normalized gate rows and avoids
weight-magnitude and precision costs altogether. The Boolean target is
a bounded-fan-in circuit; the quadratic benchmark is sufficient, not a
claim that quadratic dependence is necessary.

The theorems include actual finite mean-field priors and all stronger
scales. The same training-sample event works at every eligible learner
width and scale. Neither an infinite-width approximation nor an efficient
posterior-sampling algorithm is asserted. Earlier population-averaging
arguments motivate the model and its Gaussian estimates \cite{averaging};
all estimates needed for the present theorem are proved below.

\section{Learner, teacher, and prior restriction}

Let $d\ge1$, $L\ge2$, and $H=L-1$. Inputs lie in
$\mathcal X=\{-1,1\}^d$, with an arbitrary probability distribution $D$.
Write $\bar x=(x,1)$. A network of hidden width $N$ has
\[
 h^1=\tanh(W^1\bar x),\qquad
 h^s=\tanh(W^sh^{s-1})\quad(2\le s\le H),\qquad
 z_w(x)=v^\top h^H.
\]
There are no additional hidden biases; a constant input coordinate and
counted constant wires supply them in Boolean implementations.
The raw Gaussian law $\mu_{N,t}$ has independent centered entries with
variances
\begin{equation}\label{eq:prior}
 \frac{t^2}{d}\text{ in }W^1,\qquad
 \frac{t^2}{N}\text{ in }W^s\ (s\ge2),\qquad
 \frac{t^2}{N^2}\text{ in }v,\qquad 0<t\le1.
\end{equation}
The scale may be prescribed as a function $t=t(d,L,N)$.
In particular, $t(N)=N^{(1-\gamma)/2}$ gives every $\gamma\ge1$ in the
usual exponent family. Define the readout energy and clipped score by
\[
 E_v(w)=N\norm v_2^2,\qquad T(u)=\clip(u,-1/2,1/2).
\]

\paragraph{Population averaging.}
A balanced partition into $q$ populations has group sizes in
$\{\lfloor N/q\rfloor,\lceil N/q\rceil\}$. For each hidden layer choose
such a partition and let $P_s$ be its orthogonal averaging projection.
The averaged network $\bar w$ has
\begin{equation}\label{eq:projection}
 \bar W^1=P_1W^1,\qquad
 \bar W^s=P_sW^sP_{s-1}\ (s\ge2),\qquad
 \bar v=P_Hv.
\end{equation}
Its clean hidden activations are constant within populations. Consequently
it computes exactly the same function as a network of width $q$ obtained
by collecting equal neurons. Notice that
$E_v(\bar w)\le E_v(w)$.

\paragraph{Internal disturbances.}
For a network $u$ of width $N$, add a vector $\eta^s$ after each tanh
activation and a scalar $\eta^{\rm out}$ at the output. Denote its output
by $z_u(x;\eta)$. The disturbance norm and calibrated radius are
\begin{equation}\label{eq:noise}
 \norm\eta_*^2=\sum_{s=1}^H\frac{\norm{\eta^s}_2^2}{N}
                    +|\eta^{\rm out}|^2,
 \qquad r_c(u)=\frac{c}{\sqrt{1+H E_v(u)}},\qquad c>0.
\end{equation}
The disturbed states need not lie in $[-1,1]^N$.

\paragraph{The support condition.}
Fix $\lambda\ge1$ and $c>0$ before observing labels. Let
$\mathcal S_{N,\lambda,c}$ consist of those networks for which some
$q\le N$ and balanced partitions satisfy all three conditions:
\begin{align}
 q&\le1+E_v(w)/\lambda,\label{eq:count}\\
 \sup_{x\in\mathcal X}|T(z_w(x))-T(z_{\bar w}(x))|
   &\le1/16,\label{eq:clean}\\
 |T(z_{\bar w}(x;\eta))-T(z_{\bar w}(x))|
   &\le1/32\quad
   (x\in\mathcal X,\ \norm\eta_*\le r_c(\bar w)).\label{eq:robust}
\end{align}
In \eqref{eq:robust}, $\bar w$ denotes the averaged network in
\eqref{eq:projection}. Thus the radius is calibrated by the averaged
network's \emph{own} readout energy. The partitions may depend on the
weights, but not on the training labels.

The support is Borel. There are finitely many partition choices and inputs;
for each choice the noise supremum is over a compact ball whose radius
depends continuously on the weights. Equivalently one may rescale to a fixed
unit ball and use a countable dense set.

The learner first restricts $\mu_{N,t}$ to this support and then, for a
sample $S=((X_i,f(X_i)))_{i=1}^n$, conditions on
\begin{equation}\label{eq:training}
 f(X_i)z_w(X_i)>1/2\quad(1\le i\le n).
\end{equation}
Write the posterior as $Q_{S,N,t}$. Two further restrictions may optionally
be retained: global clean margin $|z_w(x)|>1/2$ at every input, and the
original-network internal stability condition
\begin{equation}\label{eq:optional}
 |T(z_w(x;\eta))-T(z_w(x))|\le1/16
 \quad(x\in\mathcal X,\ \norm\eta_*\le r_c(w)).
\end{equation}
The theorem holds for either, both, or neither of these optional restrictions.
When global margin is imposed, \eqref{eq:training} is equivalent to
conditioning on the exact training signs.

\paragraph{Teacher convention.}
The target is computed by a single-output, depth-$H$, width-$C$ layered
Boolean circuit with NAND, COPY, and NOT gates. A reserved constant-one
wire at every stage, including the final stage, is included in $C$, and
$C\ge2$. Compilation from a different basis must count its width and depth
overhead. A circuit of smaller logical depth can be padded with COPY gates.
Thus $L$ is the actual learner depth and the common depth after teacher
padding; it cannot be replaced by a shallower teacher's intrinsic depth
while leaving learner depth unspecified. This comparator is a Boolean
circuit, not an arbitrary real-weight neural teacher.

\section{An explicit learning theorem}

Define target constants
\begin{align}
 e_*&=\frac1{16\sqrt{2C}},&\tau_0&=e_*/4,\notag\\
 b_L&=\sqrt{16\log9+8\log(4L)},&
 A&=\max\{32b_L,2\log(8/e_*)\},\notag\\
 a&=\max\{4\sqrt\lambda,256c\},&
 \rho&=\min\{1/4,A/[4(d+1)],Ae_*/(64\sqrt C)\},\notag\\
 q_0&=d+2+(H-1)C,&
 B_0&=9dA^2+18(H-1)CA^2+2a^2C,\notag\\
 D_*&=B_0+q_0\log\frac{e\sqrt{2\pi}}{2\rho}+\log2.
 \label{eq:targetconstants}
\end{align}
For $0<\varepsilon,\delta<1$, set
\begin{align}
 B&=\max\{8,8[D_*+\log(4/\varepsilon)]\},&
 Q&=\lceil1+B/\lambda\rceil,\notag\\
 J&=8\left[D_*+\log(4/\varepsilon)+(H+1)\log Q
       +\frac{\log2}{2}\{d+2+(H-1)Q\}\right],\notag\\
 R&=1+\sqrt{QJ},&P_Q&=Q(d+2)+(H-1)Q^2,\notag\\
 \zeta&=\min\left\{\frac1{32},
 \frac{c}{2\sqrt{(1+HB)((d+1)^2+HQ^2)}}\right\}.
 \label{eq:budgets}
\end{align}

\begin{theorem}[Nearly parameter-count learning]\label{thm:main}
Let $f$ have the counted Boolean implementation above. If
\begin{equation}\label{eq:samples}
 n\ge\frac2\varepsilon\left[
 P_Q\log(3+4R/\zeta)+\log(1/\delta)\right],
\end{equation}
then, with probability at least $1-\delta$ over iid $X_i\sim D$,
\begin{equation}\label{eq:conclusion}
 \E_{w\sim Q_{S,N,t}}D\{x:f(x)z_w(x)\le1/4\}\le\varepsilon
\end{equation}
simultaneously for every integer $N\ge C$ and every $0<t\le1$.
All posterior normalizers are positive. The same conclusion holds with
any choice of the optional restrictions described above.
\end{theorem}

\begin{corollary}[Joint width--depth rate]\label{cor:rate}
Take $\lambda=d+L$. For fixed $c,\varepsilon,\delta$, the sufficient sample
bound in Theorem~\ref{thm:main} is
\[
 \widetilde O(dC+LC^2),
\]
where the soft-$O$ suppresses logarithmic factors in $d,C,L$.
More precisely, at fixed $d,c,\varepsilon,\delta$, put
\[
 F=1+\log^2(2+C)+\log(2+L),\qquad G=\log(2+CL).
\]
Then a sufficient bound is $O(C^2LF^2G)$. In particular its fixed-width
depth dependence is $O(L\log^3(2+L))$, and its fixed-depth width dependence
is $O(C^2\log^5(2+C))$.
\end{corollary}

\begin{remark}[Prediction and limiting conventions]
The posterior-average sign error is no greater than the margin error in
\eqref{eq:conclusion}. If $p_S(x)=Q_{S,N,t}\{f(x)z_w(x)>0\}$, then
$D\{p_S\le2/3\}\le3\varepsilon$ and
$D\{p_S\le1/2\}\le2\varepsilon$. Replacing $\varepsilon$ by
$\varepsilon/3$ in the budget gives strict correct-label probability above
$2/3$ outside test mass $\varepsilon$.
The theorem covers all $\gamma\ge1$ by the substitution following
\eqref{eq:prior}, and arbitrarily small positive scales at each fixed width.
It also bounds the corresponding limsup risk as $t\downarrow0$.
If parameter posteriors have a weak limit, its sign-error bound follows by
Portmanteau using the open event $fz_w<1/4$. This observation does not claim
that the closed margin event $fz_w\le1/4$ automatically transfers unchanged
at boundary mass.
\end{remark}

\section{A positive-mass robust Boolean implementation}

\begin{lemma}[Target cylinder]\label{lem:cylinder}
For every $N\ge C$ and $0<t\le1$, there is a set of weights of raw Gaussian
probability at least $\exp(-ND_*/t^2)$ which belongs to
$\mathcal S_{N,\lambda,c}$, satisfies both optional restrictions, and has
$f(x)z_w(x)>1$ at every input.
\end{lemma}
\begin{proof}
Partition the neurons into $C$ balanced clone groups of sizes $k_b$, so
$C/2\le N/k_b\le2C$. Fill unused logical slots with constants or copies.
The ideal state $b^s(x)\in\{-1,1\}^N$ is constant on each group and has
norm $\sqrt N$. Let $P_s$ average these groups.

Give every gate gain $A$. In the $\{-1,1\}$ convention NAND uses the
signed affine expression $1-u-v$, COPY uses $u$, and NOT uses $-u$.
The counted constant wire supplies the first term inside hidden layers.
Every ideal signed preactivation is at least $A$, and at most three
coefficients of magnitude $A$ enter a gate. Divide each coefficient
between two logical wires equally over the source clones. Let the readout
average the output-wire clones with total coefficient $a$. Call these
weights $w_*$.

Use standardized coordinates $\sqrt d W^1$, $\sqrt N W^s$, and $Nv$,
whose raw Gaussian variances are all $t^2$. Restrict each first-layer
coordinate, each internal row's coordinate along each of the $C$
normalized clone-group directions, and each readout coordinate to within
$\rho t$ of its teacher value. Every internal matrix can then be written
\[
 W^s=W_*^s+R^s+E^s,\qquad R^s=R^sP_{s-1},\qquad
 \norm{R^s}_{\rm op}\le\rho t\sqrt C,\qquad E^sP_{s-1}=0.
\]
The unconstrained Gaussian matrices $E^s$ are independent of these
coordinate restrictions. Two $1/4$-nets on the Euclidean unit sphere,
each of size at most $9^N$, give
\[
 \Prb\{\norm{E^s}_{\rm op}>b_Lt\}
 \le2\exp(2N\log9-Nb_L^2/8)\le\frac1{2L}.
\]
Indeed, a fixed bilinear form has variance at most $t^2/N$, and twice the
largest value over the nets bounds the operator norm. Hence all the bulk
norm bounds hold with conditional probability at least $1/2$.

We verify a stronger disturbance invariant than needed: allow RMS
magnitude at most $\tau_0$ separately at each hidden layer, and an output
disturbance at most $\tau_0$. Write
$e_s=\norm{h^s_\eta-b^s}_2/\sqrt N$. If $e_{s-1}\le e_*$, every
clone-group mean differs from its ideal bit by at most
$e_*\sqrt{N/k_b}\le1/16$. The intended signed gate preactivation
therefore remains at least $A/2$.
For $\xi\in\{-1,1\}$ the elementary inequality
\[
 \xi q\ge A/2\quad\Longrightarrow\quad
 |\tanh(q+u)-\xi|\le2e^{-A/2}+8|u|/A
\]
holds: if $|u|<A/4$, use the tanh tail bound, and otherwise the right side
is at least two. Since $E^sb^{s-1}=0$ and
$\norm{h^{s-1}_\eta}_2/\sqrt N\le1+\tau_0$, it follows that
\[
 e_s\le2e^{-A/2}+(8/A)\rho\sqrt C(1+\tau_0)
                  +(8b_L/A)e_{s-1}+\tau_0\le e_*.
\]
Each summand is at most $e_*/4$ by \eqref{eq:targetconstants}.
The first-layer preactivation perturbation is at most
$\rho(d+1)\le A/4$, which initializes the same invariant. At the readout,
\[
 f(x)z_w(x;\eta)\ge15a/16-\rho(1+\tau_0)-\tau_0>1.
\]

The readout center and perturbation satisfy
$\sqrt N\norm{v_*}_2\ge a\sqrt{C/2}$ and
$\sqrt N\norm{v-v_*}_2\le\rho$. Thus
\begin{equation}\label{eq:witnessenergy}
 E_v(w)\ge a^2C/8\ge2\lambda C,\qquad
 r_c(w)\le\frac{2\sqrt2c}{a\sqrt C}\le\tau_0.
\end{equation}
This proves original-network stability \eqref{eq:optional} with zero
clipped variation, as well as the global margin and count condition.

The clone projections map this cylinder into itself. Averaging rows
preserves the constrained intervals because teacher rows agree within each
group, and right projection kills $E^s$. First-layer and readout intervals
are preserved as well. Consequently the averaged network has the same
noisy margin invariant, and its own readout energy satisfies
\eqref{eq:witnessenergy}. Both its clean and noisy clipped scores equal
$f(x)/2$ throughout its calibrated noise ball. This proves
\eqref{eq:robust}, even with zero in place of $1/32$, and proves
\eqref{eq:clean} with zero error.

There are $Nq_0$ independent constrained coordinates; the sum of the
squared centers is at most $NB_0$. For a scalar
$Z\sim\mathcal N(0,t^2)$ and $\rho\le1/4$,
\[
 \Prb\{|Z-m|\le\rho t\}\ge
 \frac{2\rho}{e\sqrt{2\pi}}e^{-m^2/t^2}.
\]
This follows by bounding the density on the interval and using
$(|m|/t+\rho)^2/2\le m^2/t^2+\rho^2$.
Multiply these estimates and the independent bulk-event probability
$1/2$. Since $t\le1$ and $N/t^2\ge1$, the product is at least
$\exp(-ND_*/t^2)$.
\end{proof}

\section{Posterior control of adaptive population choices}

For a fixed partition tuple, let $k_{s,a}$ be its group sizes. The exact
width-$q$ quotient has coefficients
\begin{align*}
 U^1_{aj}&=k_{1,a}^{-1}\sum_{i\in a}W^1_{ij},\\
 U^s_{ab}&=k_{s,a}^{-1}\sum_{i\in a,\ j\in b}W^s_{ij},\\
 u_b&=\sum_{j\in b}v_j.
\end{align*}
Define its energy and parameter count by
\[
 \kappa_q=\frac dq\norm{U^1}_F^2+
       \sum_{s=2}^H\norm{U^s}_F^2+q\norm u_2^2,
 \qquad P_q=q(d+2)+(H-1)q^2.
\]

\begin{lemma}[A partition-uniform Gaussian tail]\label{lem:partitions}
For every integer $Q\ge1$ and $J>0$,
\begin{align*}
 &\mu_{N,t}\{\text{some balanced partition tuple with }
           q\le\min(N,Q)\text{ has }\kappa_q>J\}\\
 &\qquad\le Q\exp\left(HN\log Q+
 \frac N2[d+2+(H-1)Q]\log2-\frac{NJ}{8t^2}\right).
\end{align*}
\end{lemma}
\begin{proof}
For fixed partitions, the quotient coefficients are independent centered
Gaussians, with variances respectively
\[
 \frac{t^2}{dk_{1,a}},\qquad
 \frac{t^2k_{s-1,b}}{Nk_{s,a}},\qquad
 \frac{t^2k_{H,b}}{N^2}.
\]
After the weights in $\kappa_q$ are applied, every squared Gaussian
coordinate has variance at most $2t^2/N$. Balance gives
$N/(qk)\le2$, $qk/N\le2$, and ratios of group sizes at most two.
Therefore
\[
 \E\exp\left(\frac{N\kappa_q}{8t^2}\right)\le2^{P_q/2},
 \qquad
 \Prb(\kappa_q>J)\le e^{-NJ/(8t^2)}2^{P_q/2}.
\]
There are at most $q^{HN}$ ordered partition tuples. If $q\le N$, then
$P_q/N\le d+2+(H-1)q$. Sum the bound over
$q\le\min(N,Q)$ and use the upper bound $Q$ on that sum's length.
This controls partitions before any weight-dependent selection.
\end{proof}

\begin{lemma}[Posterior truncation]\label{lem:truncation}
Outside posterior mass at most $\varepsilon/2$, every network has
$E_v(w)\le B$ and a witnessing partition with $q\le Q$ and
$\kappa_q\le J$.
\end{lemma}
\begin{proof}
For every correctly labeled sample, Lemma~\ref{lem:cylinder} lower-bounds
the raw probability of all prior restrictions and training constraints by
$e^{-ND_*/t^2}$. Under the raw Gaussian, $NE_v/t^2$ is chi-square with
$N$ degrees of freedom. Its exponential moment gives, for $B\ge8$,
\[
 \mu_{N,t}(E_v>B)\le
 e^{-NB/(4t^2)}2^{N/2}\le e^{-NB/(8t^2)}.
\]
After dividing by the conditioning probability, the definition of $B$
bounds this posterior probability by $\varepsilon/4$. On its complement,
\eqref{eq:count} supplies $q\le Q$.

For the exceptional partition event of Lemma~\ref{lem:partitions}, use
$t\le1$ and $N\ge1$ to bound every positive exponent term by $N/t^2$
times the corresponding term in $J/8$. Its raw probability is then at most
\[
 e^{-ND_*/t^2}e^{-(N/t^2)\log(4/\varepsilon)}.
\]
Posterior normalization bounds it by $\varepsilon/4$. Off these two
exceptional events every partition of size at most $Q$ has energy at most
$J$, so any witnessing partition can be chosen.
\end{proof}

\section{Robust rounding without an exponential depth cost}

\begin{lemma}[Rounding as an internal disturbance]\label{lem:rounding}
Every retained network from Lemma~\ref{lem:truncation} has a candidate
$g$ in a fixed finite class $\mathcal G$, independent of $N,t$, such that
\[
 \sup_{x\in\mathcal X}|T(z_w(x))-T(g(x))|\le3/32,
 \qquad \log|\mathcal G|\le P_Q\log(3+4R/\zeta).
\]
\end{lemma}
\begin{proof}
Choose a witnessing partition and its quotient $(U,u)$ of width $q\le Q$.
The energy bound implies that every coefficient has magnitude at most
$\sqrt{QJ}<R$. Round each coefficient to a grid in $[-R,R]$ containing
zero, with error at most $\zeta$ and at most $3+4R/\zeta$ possible values.
Let $(U',u')$ denote the rounded quotient and $h'^s$ its clean hidden
trajectory. Tanh gives $|h'^s_a|\le1$ at every depth.

The rounded trajectory is exactly the unrounded quotient's trajectory
with the after-activation disturbances
\begin{align*}
 \eta^1_a&=\tanh((U'^1\bar x)_a)-\tanh((U^1\bar x)_a),\\
 \eta^s_a&=\tanh((U'^sh'^{s-1})_a)-\tanh((U^sh'^{s-1})_a),\\
 \eta^{\rm out}&=(u'-u)^\top h'^H.
\end{align*}
Lift each hidden disturbance constantly over its original population.
This gives the same trajectory in the averaged width-$N$ network, without
any factor depending on population sizes. By tanh's Lipschitz bound,
\[
 |\eta^1_a|\le(d+1)\zeta,\qquad
 |\eta^s_a|\le q\zeta\ (s\ge2),\qquad
 |\eta^{\rm out}|\le q\zeta.
\]
Hence the normalized total disturbance satisfies
\[
 \norm\eta_*^2\le[(d+1)^2+Hq^2]\zeta^2
 \le\frac{c^2}{4(1+HB)}\le r_c(\bar w)^2,
\]
because $E_v(\bar w)\le E_v(w)\le B$.
Applying the merged network's own robustness condition yields
\[
 \sup_x|T(z_{\bar w}(x))-T(z_{U',u'}(x))|\le1/32.
\]
Together with \eqref{eq:clean}, the total clipped approximation error is
at most $3/32$.

Zero-pad the quotient to width $Q$. This preserves its function since
$\tanh(0)=0$, and it preserves the selected grid. There are $P_Q$ physical
parameters, so the class of all width-$Q$ networks on this grid has the
claimed log cardinality. Its grids and budgets depend only on the teacher
resources, architecture, and statistical tolerances, not learner width or
Gaussian scale.
\end{proof}

\section{Generalization and the explicit rate}

\begin{proof}[Proof of Theorem~\ref{thm:main}]
Let $\psi$ be one on $(-\infty,11/32]$, zero on
$[13/32,\infty)$, and affine between these values. Any retained,
training-consistent $w$ has a candidate $g\in\mathcal G$ with zero
training losses $\psi(f(X_i)T(g(X_i)))$. Indeed, its correct clipped
training score is $1/2$ and the approximation error is at most $3/32$.
Whenever $f(x)z_w(x)\le1/4$, the candidate's correct clipped score is at
most $11/32$, so the ramp loss equals one.

For a fixed loss in $[0,1]$ of population mean $r$, the probability that
all $n$ iid losses vanish is at most $(1-r)^n\le e^{-nr}$: the probability
of a positive loss is at least its mean. A union bound over
$\mathcal G$ gives an event of probability at least $1-\delta$ on which
every zero-training-loss candidate has population loss at most
\[
 \frac{\log|\mathcal G|+\log(1/\delta)}n.
\]
This single event does not depend on $N$ or $t$. On it, integrate the
pointwise loss comparison over retained posterior networks and bound the
exceptional posterior mass by $\varepsilon/2$. Lemmas
\ref{lem:truncation} and \ref{lem:rounding} and \eqref{eq:samples} give
\[
 \E_Q D\{fz_w\le1/4\}\le\varepsilon/2+
 \frac{\log|\mathcal G|+\log(1/\delta)}n\le\varepsilon.
\]
Lemma~\ref{lem:cylinder} proves positive normalizers for every stated
choice of optional restrictions.
\end{proof}

\begin{proof}[Proof of Corollary~\ref{cor:rate}]
The target constants give
\[
 D_*=O\bigl((d+LC)\log^2(2dLC)+(\lambda+c^2)C\bigr).
\]
Thus, with $\lambda=d+L$ and fixed $c,\varepsilon,\delta$,
$Q=\widetilde O(C)$ because $d+LC\le(d+L)C$.
The quantities $J,R,\zeta^{-1}$ are polynomially bounded in $d,C,L$
up to logarithms. Consequently the grid logarithm is
$O(\log(2+d+C+L))$ up to additional logarithms of these resources, and
\eqref{eq:samples} gives $\widetilde O(dC+LC^2)$.

For the sharper fixed-$d$ statement, the target gain has $A^2=O(F)$ and
$D_*=O(CLF)$. Hence
\[
 B=O(CLF),\quad Q=O(CF),\quad J=O(CLF),\quad
 R=O(C\sqrt L F),\quad P_Q=O(C^2LF^2).
\]
The explicit mesh has a polynomial inverse in $C,L,F$; for example
$\zeta^{-1}=O(1+C^{3/2}L^{3/2}F^{3/2})$ at fixed $d,c$.
Therefore $\log(3+4R/\zeta)=O(G)$. Multiplication gives
$O(C^2LF^2G)$ and the stated specializations.
\end{proof}

\section{Removing depth logarithms with per-layer stability}
\label{sec:perlayer}

A stronger disturbance condition removes even the depth logarithms from
Corollary~\ref{cor:rate}, while continuing to allow analog hidden states.
The condition is a product of noise balls: every layer may be perturbed
by its full allowed amount simultaneously. We also use a Euclidean-ball
parameter count, so that merely having more layers does not enlarge the
required range of a typical coded coefficient.

Fix $\lambda\ge1$ and $c>0$. Retain the balanced-population count
\eqref{eq:count} and clean approximation \eqref{eq:clean}, but replace
\eqref{eq:robust} by
\begin{equation}\label{eq:perlayer-stability}
 |T(z_{\bar w}(x;\eta))-T(z_{\bar w}(x))|\le\frac1{32}
 \quad\text{whenever}\quad
 \frac{\norm{\eta^s}_2}{\sqrt N}\le\tau_c(\bar w)\ (1\le s\le H),
 \quad |\eta^{\rm out}|\le\tau_c(\bar w),
\end{equation}
where
\[
 \tau_c(\bar w)=\frac{c}{\sqrt{1+E_v(\bar w)/\lambda}}.
\]
This uses the averaged full-width network's own readout energy. It is a
stronger disturbance requirement than the aggregate-radius condition:
any vector of aggregate norm at most $r_c(\bar w)$ obeys each bound above,
since $H\ge1$ and $\lambda\ge1$. As before, the prior can optionally
retain global margin and original-network stability \eqref{eq:optional}.
The training condition remains \eqref{eq:training}.

For this support, replace the target constants in
\eqref{eq:targetconstants} by
\begin{align}
 b&=\sqrt{16\log9+8\log4},&
 A&=\max\{32b,2\log(8/e_*)\},\notag\\
 a&=\sqrt\lambda\max\{4,256c\},&
 D_{\rm pl}&=B_0+q_0\log\frac{e\sqrt{2\pi}}{2\rho}+(H-1)\log2.
 \label{eq:perlayer-target}
\end{align}
Here $e_*,\tau_0,\rho,q_0,B_0$ retain their previous formulas, evaluated
with the new $A,a$. In particular the gate gain $A$ no longer depends on
depth. Define $B,Q,J,P_Q$ by \eqref{eq:budgets}, with $D_{\rm pl}$ in
place of $D_*$, and set
\begin{equation}\label{eq:perlayer-mesh}
 \tau_B=\frac c{\sqrt{1+B/\lambda}},\qquad
 \zeta_{\rm pl}=\frac{\tau_B}{2\max\{d+1,Q\}},\qquad
 K_{\rm ball}=\sqrt{2\pi e}.
\end{equation}

\begin{theorem}[Linear depth under per-layer stability]
\label{thm:perlayer}
Under the preceding support and teacher conventions, the conclusions of
Theorem~\ref{thm:main} hold if
\begin{equation}\label{eq:perlayer-sample}
 n\ge\frac2\varepsilon\left[
 P_Q\log\left(K_{\rm ball}
       \left(1+\frac{\sqrt{QJ/P_Q}}{\zeta_{\rm pl}}\right)\right)
       +\log(1/\delta)\right].
\end{equation}
For $\lambda=d+L$ and fixed $d,c,\varepsilon,\delta$, this sufficient
bound is
\[
 O\bigl(C^2L\log^5(2+C)\bigr).
\]
In particular, there is no hidden logarithm of depth in this rate.
\end{theorem}

\begin{proof}
We first check the target cylinder. At each internal layer, the same
Gaussian two-net bound as before now gives
\[
 \Prb\{\norm{E^s}_{\rm op}>bt\}
 \le2\exp(2N\log9-Nb^2/8)=2\cdot4^{-N}\le\frac12.
\]
The internal bulk matrices are independent of one another and of the
constrained scalar coordinates. Thus all $H-1$ bulk bounds hold with
probability at least $2^{-(H-1)}$. Paying this probability directly,
rather than choosing a cutoff that makes their intersection probability
at least $1/2$, removes the depth dependence from $b$ and $A$.
The scalar interval calculation in Lemma~\ref{lem:cylinder} gives raw
cylinder probability at least $e^{-ND_{\rm pl}/t^2}$.

The same saturation induction proves signed output greater than one under
RMS perturbation at most $\tau_0$ separately at every hidden layer and
output. Projection maps the cylinder into itself, so both $w$ and
$\bar w$ satisfy this invariant. For the averaged network,
\[
 E_v(\bar w)\ge a^2C/8,\qquad
 \tau_c(\bar w)\le\frac{c\sqrt{8\lambda}}{a\sqrt C}
                  \le\tau_0.
\]
The new per-layer condition holds with zero clipped variation. Also
$E_v(w)\ge2\lambda C$, proving the population-count condition, and
$a\ge256c$ proves the optional original-network stability as before.
The target normalizer bound is therefore valid for each stated prior.

Lemma~\ref{lem:truncation} applies unchanged with this target cost.
Outside posterior mass $\varepsilon/2$, there is a witnessing quotient
of width $q\le Q$, with $\kappa_q\le J$, and
$E_v(\bar w)\le E_v(w)\le B$. Hence its permissible per-layer radius
is at least $\tau_B$.

Round quotient coefficients to the nearest multiple of
$\zeta_{\rm pl}$. The construction in Lemma~\ref{lem:rounding}
expresses the rounded trajectory exactly through population-constant
internal disturbances. Every first-layer coordinate disturbance is at
most $(d+1)\zeta_{\rm pl}$, and every subsequent hidden or output
disturbance is at most $q\zeta_{\rm pl}$. Consequently each layer's RMS
norm and the output disturbance are at most $\tau_B/2$. The per-layer
condition bounds the rounded-to-averaged clipped error by $1/32$, and
the total approximation error is again at most $3/32$.

It remains to count rounded coefficients without a depth-dependent range
penalty. Since $d\ge1$ and $q\le Q$, the sum of squares of all unweighted
quotient parameters is at most $QJ$. After zero-padding to width $Q$, their
vector lies in a Euclidean ball in dimension $P_Q$ of radius
$S=\sqrt{QJ}$. Consider all lattice points produced by nearest-grid
rounding of any vector in this ball. Their cubes of side
$\zeta_{\rm pl}$ have disjoint interiors and lie in the ball of radius
$S+\zeta_{\rm pl}\sqrt{P_Q}$.
For every positive integer $P$ and radius $R$,
\[
 \mathrm{vol}(B_R)\le
 \left(\sqrt{\frac{2\pi e}{P}}R\right)^P.
\]
For example, integrate
$\mathbf1_{\norm x\le R}\le\exp(\alpha R^2-\alpha\norm x^2)$ and
choose $\alpha=P/(2R^2)$. Dividing by the lattice-cube volume yields
\[
 |\mathcal G|\le
 \left[K_{\rm ball}
       \left(1+\frac{\sqrt{QJ/P_Q}}{\zeta_{\rm pl}}\right)\right]^{P_Q}.
\]
Padding includes every $q\le Q$ in the same ball, so no separate count of
partitions or widths is needed at this step. The grid candidates need
not themselves satisfy the prior restrictions.

The ramp and finite-class argument from Theorem~\ref{thm:main} now give
\eqref{eq:perlayer-sample}, on a single event for every $N\ge C$ and
$0<t\le1$.

For the rate, put $F_C=1+\log^2(2+C)$ and fix $d,c,\varepsilon,\delta$.
At $\lambda=d+L$, the modified target cylinder has
\[
 D_{\rm pl}=O(LCF_C),\qquad B=O(LCF_C),\qquad
 Q=O(CF_C),\qquad J=O(LCF_C).
\]
More specifically, $D_{\rm pl}+\log(4/\varepsilon)\le B/8$ and
$B/\lambda\le Q$. The formula for $J$ therefore gives $J=O(LQ)$ at
fixed dimension. If $H\ge2$, then
$P_Q\ge(H-1)Q^2$, so $QJ/P_Q=O(1)$ uniformly in depth. Also
\[
 \zeta_{\rm pl}^{-1}
 =\frac{2\max\{d+1,Q\}\sqrt{1+B/\lambda}}c
 =O(Q^{3/2})
\]
at fixed $d,c$. Thus the logarithm in \eqref{eq:perlayer-sample} is
$O(\log(2+C))$, while $P_Q=O(LC^2F_C^2)$. If $H=1$, then $L=2$
is fixed and the same claimed bound follows directly from the displayed
budgets. This proves $O(C^2L\log^5(2+C))$.
\end{proof}

The two changes have distinct roles. Requiring simultaneous per-layer
robustness makes the rounding precision independent of depth, and counting
an energy ball rather than a coordinate cube removes a depth factor from
the apparent coefficient range. Neither follows merely from clean
population averaging. This theorem still permits analog population
responses; the digital-margin condition introduced separately gives a
stronger, literal $C^2L$ bound.


\section{A literal quadratic-width, linear-depth theorem}
\label{sec:digital}

The logarithms in the preceding parameter-grid bound can also be removed
under a different structural restriction: after merging, the network behaves as a
collection of robust digital gates.  This assumption is natural for a
saturated implementation of a Boolean circuit, but is not claimed to hold
for typical gradient-trained analog networks.  The resulting prior still
has positive probability around an implementation of every eligible
Boolean circuit.  Its support definition does not use that circuit or its
width.

We retain the Gaussian tanh model, $H=L-1$, the readout energy
$E_v(w)=N\|v\|_2^2$, and the clipped score
$T(u)=\max\{-1/2,\min\{u,1/2\}\}$.
For balanced population partitions into $q$ groups at every hidden layer,
let $\bar w$ be the population-averaged full-width network.  Write its
exactly equivalent width-$q$ quotient as
\[
 h^0(x)=(x,1),\qquad a^s(x)=U^sh^{s-1}(x),\qquad
 h^s(x)=\tanh(a^s(x)),\qquad z_{\bar w}(x)=u^\top h^H(x).
\]
Thus $U^1$ has $d+1$ columns, subsequent $U^s$ have $q$ columns,
and $u$ has $q$ coordinates.

Fix, before observing labels,
\begin{equation}\label{eq:digital-calibration}
 \lambda=(d+L)\bigl(1+\log(d+2)\bigr),\qquad
 \Psi(q)=q\log^2(2+q),\qquad
 \eta=\frac1{16},\quad \kappa=\frac15.
\end{equation}
The event $\mathcal D_N$ consists of networks admitting such partitions,
with $1\le q\le N$, for which
\begin{align}
 \lambda\Psi(q)&\le E_v(w),\label{eq:digital-count}\\
 \sup_x|T(z_w(x))-T(z_{\bar w}(x))|&\le\frac1{16},
                    \label{eq:digital-merge}\\
 |h^s_a(x)|&\ge\frac{15}{16},
                    \label{eq:digital-saturation}\\
 |a^s_a(x)|&\ge\kappa\|U^s_{a,\cdot}\|_1,
 \qquad |z_{\bar w}(x)|\ge\kappa\|u\|_1.
                    \label{eq:digital-relative}
\end{align}
The last two lines hold for every Boolean input, every hidden layer,
and every indicated coordinate.  All rows and the readout in
\eqref{eq:digital-relative} are required to be nonzero.
The prior is the raw Gaussian conditioned on $\mathcal D_N$ and
\[
 \mathcal G_N=\{w: |z_w(x)|>1/2\text{ for every Boolean }x\}.
\]
Training conditions further on the observed noiseless signs.  We also
allow intersection with the original internal-noise event
$\mathcal R_{N,c}$, for any fixed $c>0$: clipped scores must change by at
most $1/16$ under every simultaneous post-activation and output disturbance
whose norm
\[
 \left(\sum_{s=1}^H\|\xi^s\|_2^2/N+|\xi^{\rm out}|^2\right)^{1/2}
\]
is at most $c/\sqrt{1+HE_v(w)}$.
This optional intersection does not change the conclusion below.

These are Borel restrictions involving finitely many inputs and partition
tuples at each width.  Saturation makes the populations nearly Boolean;
the relative row margin asks that the sign of each gate tolerate a fixed
fractional error in its aggregate incoming coefficients.  The stronger
population-count calibration in \eqref{eq:digital-count} is an additional
choice of prior: it compensates for the logarithmic gains used in the
Boolean implementation.  It is not implied by the original robustness
restriction.

\subsection{Exact budgets and the learning statement}

Let the Boolean target have layered width $C$, with $H$ stages and the
constant wires required by the NAND/COPY/NOT implementation counted.
Put
\begin{align}
 e_*&=\frac1{16\sqrt{2C}},\qquad \tau_0=e_*/4,\qquad
 b=\sqrt{16\log9+8\log4},\notag\\
 A&=\max\{32b,2\log(8/e_*)\},\qquad
 a=\sqrt\lambda\max\{4\log(2+C),256c\},\notag\\
 \rho&=\min\{1/4,A/[4(d+1)],Ae_*/(64\sqrt C)\},\notag\\
 q_0&=d+2+(H-1)C,\notag\\
 B_0&=9dA^2+18(H-1)CA^2+2a^2C,\notag\\
 D_{\rm dig}&=B_0+q_0\log\frac{e\sqrt{2\pi}}{2\rho}+(H-1)\log2.
 \label{eq:digital-target-budget}
\end{align}
For $0<\varepsilon,\delta<1$, define
\begin{equation}\label{eq:digital-stat-budget}
 \begin{aligned}
 B&=\max\{8,8[D_{\rm dig}+\log(6/\varepsilon)]\},\\
 Q&=\max\{q\in\mathbb N:q\ge1,\ \Psi(q)\le B/\lambda\},\\
 P_Q&=Q(d+2)+(H-1)Q^2.
 \end{aligned}
\end{equation}
The maximum exists and satisfies $Q\ge C$, as the proof verifies.

\begin{theorem}[Literal quadratic width and linear depth]
\label{thm:digital}
Under either prior just specified, all posterior normalizers are positive.
If
\begin{equation}\label{eq:digital-samples}
 n\ge\frac6\varepsilon
       \left[\log Q+P_Q\log33+\log(1/\delta)\right],
\end{equation}
then, with probability at least $1-\delta$ over iid training inputs,
the posterior gives the correct label probability strictly greater than
$2/3$ outside a set of test inputs of probability at most $\varepsilon$.
The same sample event works for every $N\ge C$ and every $0<t\le1$,
under an arbitrary common training and test distribution on the Boolean
cube.

There is an absolute constant $K$ such that the simpler sufficient bound
\begin{equation}\label{eq:digital-envelope}
 n\ge\frac{K(1+c^2)^2}{\varepsilon}
       \left[(d+2)R+LR^2+\log(1/\delta)\right],
 \qquad R=C+\frac{\log(6/\varepsilon)}{\lambda},
\end{equation}
implies the same conclusion.  In particular, at fixed input dimension,
calibration, accuracy, and confidence, the sufficient sample count is
$O(C^2L)$, with no hidden logarithm of $C$ or $L$.
\end{theorem}

The theorem covers the finite mean-field prior and every stronger scaling
$\gamma\ge1$, by taking $t(N)=N^{(1-\gamma)/2}$.  Its digital hypothesis is separate from ordinary internal-noise stability;
the conclusion also holds when both conditions are imposed. It strengthens
clean population averaging through the stated digital and count requirements.  It gives no claim about efficient posterior
sampling or typical training outcomes.

\subsection{A Gaussian cylinder inside the digital support}

\begin{lemma}[Digital target cylinder]\label{lem:digital-cylinder}
For every $N\ge C$ and $0<t\le1$, a set of raw probability at least
$\exp(-ND_{\rm dig}/t^2)$ belongs to
$\mathcal G_N\cap\mathcal D_N\cap\mathcal R_{N,c}$ and satisfies
$f(x)z_w(x)>1$ at every input.
\end{lemma}
\begin{proof}
Use the balanced cloned Boolean implementation with gate gain $A$ and
readout coefficient sum $a$.  In coordinates
$\sqrt d W^1$, $\sqrt N W^s$, and $Nv$, constrain first-layer coordinates,
every internal row's normalized clone-group directions, and all readout
coordinates to intervals of radius $\rho t$ around their teacher values.
The unconstrained internal blocks satisfy $E^sP=0$.  The same two-net
Gaussian estimate used for the target cylinder gives
\[
 \Pr\{\|E^s\|_{\rm op}>bt\}
 \le2\exp(2N\log9-Nb^2/8)=2\cdot4^{-N}\le\frac12.
\]
The blocks are independent.  Requiring all $H-1$ bulk bounds therefore
costs at most a factor $2^{-(H-1)}$.  This replaces a depth-dependent
operator-norm cutoff by a constant cutoff; the probability cost is
included in $D_{\rm dig}$.

The saturation induction for the cloned construction gives hidden RMS
error at most $e_*$ relative to its ideal Boolean state, even with RMS
noise at most $\tau_0$ separately at every layer and output noise at most
$\tau_0$.  Indeed its recurrence is
\[
 e_s\le 2e^{-A/2}+(8/A)\rho\sqrt C(1+\tau_0)
                 +(8b/A)e_{s-1}+\tau_0\le e_*.
\]
The first-layer perturbation initializes the same invariant, and the
signed readout is at least
$15a/16-\rho(1+\tau_0)-\tau_0>1$.
There are $Nq_0$ independent constrained scalar coordinates with sum of
squared centers at most $NB_0$.  The scalar Gaussian interval estimate
\[
 \Pr\{|Z-m|\le\rho t\}
 \ge\frac{2\rho}{e\sqrt{2\pi}}e^{-m^2/t^2}
 \quad (Z\sim\mathcal N(0,t^2))
\]
and $N/t^2\ge1$ give the claimed probability $e^{-ND_{\rm dig}/t^2}$.

Averaging within clone groups maps this cylinder into itself.  Both the
original and averaged networks therefore have clipped score $f/2$,
proving \eqref{eq:digital-merge} and the global margin.
Furthermore,
\[
 E_v(w)\ge a^2C/8\ge2\lambda C\log^2(2+C),
\]
which proves \eqref{eq:digital-count}.  The calibration $a\ge256c\sqrt\lambda$
also ensures that the original internal-noise ball lies in the preceding
$\tau_0$ product of balls, establishing the optional restriction.

It remains to check the digital conditions for the quotient.  Its hidden
coordinates are constant on groups of mass at least $1/(2C)$.  The RMS
error bound thus implies that every coordinate differs from its ideal
bit by at most $\sqrt{2C}e_*=1/16$, proving saturation.
Every ideal gate row has $\ell_1$ norm at most $3A$ and signed ideal
preactivation at least $A$.

In the first layer, the quotient row perturbation has $\ell_1$ norm at
most $(d+1)\rho t/\sqrt d\le A/4$.  Signed preactivation divided by row
norm is therefore at least $3/13>1/5$.
For an internal row, the perturbation coefficient from group $b$ is at
most $\rho t\sqrt{k_b/N}$.  Its total $\ell_1$ norm is at most
\[
 \rho\sqrt{2C}\le A/1024.
\]
Accounting also for the coordinatewise $1/16$ error of the preceding
state gives signed preactivation at least $831A/1024$ and row norm at
most $3073A/1024$.  Their ratio exceeds $1/5$.
Finally, the quotient readout perturbation has $\ell_1$ norm at most
$\rho$, so its signed output is at least $15a/16-\rho$ and its row norm
at most $a+\rho$.  As $a\ge4$ and $\rho\le1/4$, their ratio is at least
$14/17>1/5$.  All these rows are nonzero.  This verifies
\eqref{eq:digital-relative} and completes the support check.
\end{proof}

\subsection{A constant-precision digital code}

\begin{lemma}[Finite digital class]\label{lem:digital-code}
For each integer $Q\ge1$, the sign functions of networks in
$\mathcal G_N\cap\mathcal D_N$ admitting witnesses with $q\le Q$, over
all learner widths $N$, belong to a fixed finite class $\mathcal F_Q$
with
\[
 \log|\mathcal F_Q|\le\log Q+P_Q\log33.
\]
This class does not depend on the sample, learner width, or Gaussian scale.
\end{lemma}
\begin{proof}
The unit $\ell_1$ ball in $\mathbb R^m$ has an $\ell_1$ net of radius
$1/16$ with at most $33^m$ points.  To see this, take a maximal
$1/16$-separated set.  Disjoint balls of radius $1/32$ around its points
fit inside the ball of radius $1+1/32$; volume scaling gives the bound,
and maximality gives a net.

For a quotient satisfying the digital conditions, let
$b^s(x)=\operatorname{sign}(h^s(x))$.  Normalize each nonzero row by its
$\ell_1$ norm.  Since $\|h^s-b^s\|_\infty\le1/16$, replacing a preceding
hidden state by its bits changes a normalized preactivation by at most
$1/16$.  Each hidden gate and the final readout on those bits consequently
have the original sign with margin at least
\[
 \kappa-\eta=\frac{11}{80}.
\]
The first-layer input is already Boolean, including its constant
coordinate, so it has margin at least $\kappa$.

Replace every normalized row by a point in the fixed $1/16$ net for its
dimension, and use threshold gates instead of tanh.  On Boolean inputs
this changes each preactivation by at most $1/16<11/80$.  Induction through
the layers preserves every hidden sign and the readout sign.  At width
$q$, the product of the row-net sizes is at most $33^{P_q}$.  Summing over
$q\le Q$ gives at most $Q33^{P_Q}$ possibilities.  No count of adaptive
population partitions is required: every resulting quotient is already
covered by these same row nets.

Finally, the original global margin and
\eqref{eq:digital-merge} force the original and quotient scores to have
the same sign at every input.  Thus this code describes the actual
posterior sign functions, not just an approximate surrogate.
\end{proof}

The finite nets may use real coefficients.  The lemma counts a fixed
finite function class; it does not claim a rational bit encoder or an
efficient construction algorithm.

\subsection{Posterior learning and the explicit polynomial}

\begin{proof}[Proof of Theorem~\ref{thm:digital}]
Lemma~\ref{lem:digital-cylinder} lower-bounds the raw conditioning
probability by $e^{-ND_{\rm dig}/t^2}$ for every correctly labeled sample.  Under
the raw prior, $NE_v/t^2$ is chi-square with $N$ degrees of freedom.
For $B\ge8$, its exponential moment gives
\[
 \mu_{N,t}\{E_v>B\}\le e^{-NB/(8t^2)}.
\]
After conditioning on prior support and training labels, the mass of
this event is therefore at most
\[
 \exp\left[-\frac N{t^2}(B/8-D_{\rm dig})\right]\le\varepsilon/6.
\]
Every remaining predictor has a witness with $q\le Q$.  The definition
of $Q$ is legitimate because
$D_{\rm dig}\ge2a^2C\ge32\lambda\Psi(C)$, so $B/\lambda\ge256\Psi(C)$.

Every retained sign function lies in $\mathcal F_Q$ and is consistent
with the sample.  A fixed function of population error greater than
$\varepsilon/6$ is consistent with $n$ iid labels with probability at
most $e^{-n\varepsilon/6}$.  Lemma~\ref{lem:digital-code}, a union bound,
and \eqref{eq:digital-samples} imply that, with probability at least
$1-\delta$, every consistent member of $\mathcal F_Q$ has error at most
$\varepsilon/6$.  This event is independent of $N,t$.
On it, the full posterior Gibbs error is at most $\varepsilon/3$:
$\varepsilon/6$ for retained mass and at most $\varepsilon/6$ for
discarded mass.  The test mass where correct-label posterior probability
is at most $2/3$ is at most three times this Gibbs error.  This proves the
stated confidence guarantee, including arbitrary posterior mixtures.

For the simpler envelope, all constants below are absolute.  From
\eqref{eq:digital-target-budget},
\[
 A^2\le K_1\log^2(2+C),\qquad
 \rho^{-1}\le K_2(d+C+1),\qquad
 a^2\le K_3(1+c^2)\lambda\log^2(2+C).
\]
Using $d+LC\le(d+L)C$ and the dimension logarithm in
\eqref{eq:digital-calibration}, these inequalities give
\begin{equation}\label{eq:digital-cost-envelope}
 D_{\rm dig}\le K_4(1+c^2)\lambda\Psi(C).
\end{equation}
The interval logarithm contributes only
$O(\log^2(2+C)+\log(d+2))$, and the bulk cost is linear in $H$;
there is no logarithm of $L$ in this estimate.

For $k\ge1$, $\Psi(kC)\ge k\Psi(C)$.  Moreover, for $q\ge kC$,
\[
 \Psi(q)\ge k\Psi(C)+(q-kC)\log^2 3.
\]
Combining this with \eqref{eq:digital-cost-envelope} and the definition
of $B$ bounds
\[
 Q\le K_5\left((1+c^2)C+
                  \frac{\log(6/\varepsilon)}\lambda\right)
       \le K_5(1+c^2)R.
\]
Since $\log Q\le Q$ and
$P_Q\le(d+2)Q+LQ^2$, substitution in
\eqref{eq:digital-samples} proves \eqref{eq:digital-envelope} after
enlarging one absolute constant.
\end{proof}

\begin{corollary}[A polynomial accuracy envelope]
\label{cor:digital-polynomial}
For an absolute constant $K$, the sufficient sample condition
\begin{equation}\label{eq:digital-clean-envelope}
 n\ge\frac K\varepsilon\left[
 (1+c^2)^2(dC+LC^2)+\log^2(6/\varepsilon)+\log(1/\delta)\right]
\end{equation}
implies every conclusion of Theorem~\ref{thm:digital}.
\end{corollary}
\begin{proof}
Write $s=1+c^2$ and $v=\log(6/\varepsilon)>1$.  The preceding proof gives
$Q\le K_5(sC+v/\lambda)$.  Since $\lambda\ge d+L$,
\[
 (d+2)\frac v\lambda=O(v),\qquad
 L\frac{v^2}{\lambda^2}\le v^2,\qquad
 2LsC\frac v\lambda\le2sCv\le s^2C^2+v^2.
\]
Also $(d+2)sC=O(s^2dC)$, because $d\ge1$ and $s\ge1$.
Consequently
$\log Q+P_Q\log33=O(s^2(dC+LC^2)+v^2)$.
Insert this estimate into \eqref{eq:digital-samples}.
\end{proof}

\paragraph{What the sharpening does and does not establish.}
The $C^2L$ conclusion uses both digital relative margins and the stronger
readout-energy population calibration.  It does not assert the same bound
for the original internal-noise prior, for clean neuron averaging alone,
or for the unrestricted Gaussian prior.  Every eligible Boolean circuit
has a positive-mass witness, but that is an existence statement about an
implementation, not evidence that generic trained networks satisfy the
support condition.  The proof is a conventional argument; no Lean
formalization is claimed for this theorem.


\section{Interpretation and limitations}

The analog bounds are close, up to width logarithms, to the dense parameter
count of a width-$C$, depth-$L$ representation. The digital theorem removes
those logarithms under a different structural condition and calibration.
These are upper bounds from specified priors and proofs, not optimality
theorems. A bounded-fan-in
Boolean circuit can have a description of only $O(LC\log(C+d))$ bits, so
quadratic width is not a universal lower bound for the target class.

The analog stability hypothesis has an operational meaning: a compressed or merged
network must remain stable after it replaces the original model. Matching
clean outputs alone does not ensure reliable behavior under numerical or
internal perturbations. The Boolean witnesses have this property because
saturated gates repeatedly restore approximate bits to their intended
logical values. The support condition nevertheless remains strong. No
claim is made that SGD or Adam typically produces it, or that an efficient
algorithm finds its witnessing populations.

The original network's internal robustness alone does not supply the new
argument. Clean-output closeness of two networks does not compare their
noisy hidden trajectories. Requiring robustness of the averaged
representation closes exactly this gap. The architecture-dependent count
calibration $\lambda=d+L$ is also part of the changed prior; it is not a
free improvement for the baseline population count. The digital theorem
instead uses $\lambda=(d+L)(1+\log(d+2))$ and a logarithmically weighted
population budget; its global output margin is needed to transfer the
finite sign code back to the actual posterior predictor.

The proof uses bounded hidden activations at the rounding step. Extending
this exact argument to an unbounded activation requires a suitable bound
on the rounded trajectory or a separate control of large neuron activity.
The theorem as stated is for tanh, and the target-mass proof uses its
bit-restoring saturation. Finally, the result guarantees statistical
performance of ideal Bayesian conditioning and does not provide an
efficient inference algorithm.

\paragraph{Verification status.}
This manuscript gives conventional proofs and has not been formalized in
Lean. Its main ingredients are adapted and sharpened from the population
averaging analysis \cite{averaging}; prior formal verification of related
learning theorems does not certify this new support condition or its
sample bound.

\begin{thebibliography}{9}
\bibitem{averaging}
ChatGPT-6 Astra, Claude Fable 5 and Opus 5, Kaarel H\"anni, and Dmitry Vaintrob.
\emph{Quadratic-width learning bounds for Bayesian networks with stable
neuron populations}. Research draft, 21 October 2026.
\href{https://github.com/kaarelh/AI-works/tree/main/neural-network-generalization/works-in-progress/improved-polynomials/neuron-averaging}{Repository manuscript and source}.
\end{thebibliography}
\end{document}
